% P10-4a, 14.09.2026: Mechanismus, Quellenreichweite und v4-Nachweis angeglichen.
% M08, 11.09.2026: Speichermechanismus, Projektion und Rückweg getrennt.
% Geprüft: JsonCoreRepository, CoreKangal, GithubIssueTemplate,
% GithubDriftCheck, GithubSnapshotGuard, GithubInboundHarvest,
% PipelineFullRunner (GitHub-Eingang), GithubDocPublish.
\section{Projektseitige Kernmechanismen: Persistenz und Außenkopplung}
\label{sec:realisierung-persistenz-aussenkopplung}

Die in den Abschnitten~\ref{sec:core-projektfortschreibung}
und~\ref{sec:projektionen-kontrollierte-aussenkopplung} beschriebenen
Mechanismen liegen an den Grenzen der MAF-Orchestrierung.
Das Framework verbindet die Verarbeitungsschritte; das fachliche
Zustandsmodell, seine Speicherregeln und die Außenabbildung sind
projektspezifisch realisiert.

\emph{Speicherung.} Die gemeinsame Speichernaht ist als
\texttt{SaveAsync} des \texttt{JsonCoreRepository} umgesetzt.
Vor dem Schreiben prüft \texttt{CoreKangal} strukturelle Regeln,
etwa die Endpunkte interner Beziehungen und die eindeutige
Issue-Zuordnung. Die Existenz externer Claim-Ziele und die
vollständige Auflösung des Herkunftsregisters werden hier nicht
geprüft. Fehler brechen den Speichervorgang ab, Warnungen
werden protokolliert. Erst danach wird bei verändertem Dateiinhalt
der bisherige Stand als Sicherungskopie (Snapshot) gesichert und die neue Fassung
geschrieben. Die Historienkopie setzt also einen vorhandenen und
tatsächlich abweichenden Vorgänger voraus.
Listing~\ref{lst:speichernaht} zeigt die Reihenfolge.

% M12: Das kurze Listing mit seiner Beschriftung zusammenhalten.
\par\noindent
\begin{minipage}{\linewidth}
% M12-Nachtrag: Darstellung nur fuer dieses Listing, nicht in der Uni-Vorlage.
\lstset{
  language=[Sharp]C,
  basicstyle=\ttfamily,
  columns=fullflexible,
  keepspaces=true,
  showstringspaces=false,
  breaklines=true,
  breakatwhitespace=true
}
\begin{lstlisting}[caption={[Speicherung nach Integritätsprüfung]Speicherung nach Integritätsprüfung
  (Ausschnitt aus \texttt{JsonCoreRepository.cs}).},
  label={lst:speichernaht}]
var watch = CoreKangal.Check(core);
foreach (var w in watch.Warnings)
    Console.Error.WriteLine(
        $"[core-kangal] WARNUNG {w.Code}: ...");
if (!watch.Pass)
    throw new InvalidOperationException(
        "CORE_KANGAL: Save abgebrochen - ...");

var json = JsonSerializer.Serialize(
    core, ProjectStateJson.Options);
await SnapshotBeforeOverwriteAsync(json, ct)
    .ConfigureAwait(false);
await File.WriteAllTextAsync(_file, json, ct)
    .ConfigureAwait(false);
\end{lstlisting}
\end{minipage}
\par\medskip

Diese Operation prüft die Struktur des übergebenen Bestands,
autorisiert aber keine Änderung. Die menschliche Auswahl erfolgt
zuvor in den fachlichen Verarbeitungspfaden. Beim Besuchsfall F1
liegt sie beispielsweise in den acht übernommenen und vier
abgewiesenen Operationen der Anforderungseinordnung; eine erfolgreiche Speicherprüfung
ersetzt diese Auswahl nicht. Die Prüfung der im Code identifizierten
Schreibpfade stützt ihre Zuordnung zur gemeinsamen Speicheroperation;
sie schließt nicht jede denkbare Umgehung aus. Der initiale Seed bleibt
die in Abschnitt~\ref{sec:endtoend-pfade} erklärte Ausnahme von
der Einzelautorisierung. Abschnitt~\ref{sec:eval-kontrollen}
untersucht ausgewählte Zustandswirkungen; ein Angriffstest außerhalb
der Speichernaht fehlt.

\emph{Issue-Darstellung und Änderungsvergleich.}
\texttt{GithubIssueTemplate} bildet die vorgesehenen Core-Inhalte
ohne Modellaufruf in die Issue-Struktur ab. Die Beziehung zwischen
PBI und Issue wird im Core geführt und strukturell geprüft.
Um externe Änderungen von neuer Core-Information zu unterscheiden,
vergleicht \texttt{GithubDriftCheck} Titel und Inhalt des erhobenen
Issues mit den beim letzten eigenen Schreiben gespeicherten
Prüfwerten (Hashwerten). Der Vergleich erfolgt bewusst nicht gegen eine neu
gerenderte Core-Fassung. Normalisiert werden unter anderem
Zeilenenden und nachgestellte Leerzeichen. Fehlen die gespeicherten
Vergleichswerte, ist die Abweichung nicht bestimmbar; dies ist kein
Nachweis eines unveränderten Issues.

\emph{Schnappschüsse und Rückweg.} Die Prüfung durch
\texttt{GithubSnapshotGuard} vergleicht bei konfiguriertem Ziel
den Repository-Stempel mit diesem Ziel und weist fremde oder
ungestempelte Bestände zurück. Sie ist in die geprüften
Einstiege für die GitHub-Veröffentlichung (Forward) eingebunden
und prüft keine Aktualität.
Davon getrennt ruft der reguläre GitHub-Gesamteingang Issues und
Kommentare neu ab. \texttt{GithubInboundHarvest} weist zusätzlich
Issue-Schnappschüsse zurück, deren ermitteltes Alter 24 Stunden
überschreitet. Diese Altersgrenze garantiert weder einen lückenlosen
Außenbestand noch einen weiterhin aktuellen Stand beim späteren
Schreiben. Für manuell übergebene Schnappschüsse dieses Rückwegs
ist kein Aufruf der genannten Repository-Zugehörigkeitsprüfung
eingebunden.

Sammlung und Erkennung der Rückmeldungen erfolgen regelbasiert.
Bei verarbeitbaren Issue- oder Kommentarfunden werden hingegen
MAF-Agenten zur fachlichen Aufbereitung aufgerufen; ohne solche
Funde entfällt der Modellaufruf. Ihre Vorschläge werden in ein
Delta überführt und dem regulären Einordnungs- und
Autorisierungspfad übergeben. Fehlt ein verarbeitbares Delta,
endet der Gesamteingang vor diesem Pfad. Dokumentveröffentlichungen
nutzen dieselbe externe Freigabeinstanz wie Issues, besitzen aber
keinen entsprechenden Rückweg
(Abschnitt~\ref{sec:projektionen-kontrollierte-aussenkopplung}).

Die technische Trennung folgt damit der Architekturbegründung:
Modelle deuten neue Inhalte, explizite Regeln sichern Speicherung
und Abbildung bereits festgelegter Inhalte. Wie die Ausführung
beobachtbar gemacht wird, behandelt der folgende Abschnitt.
